% Fuente: Auxiliar 1, P4 (pauta).
El objetivo es minimizar el costo de producción y transporte. Los subíndices serán:
\begin{itemize}
    \item $i$: planta
    \item $j$: zona
    \item $k$: producto
\end{itemize}

Las variables de decisión del modelo serán:
\begin{itemize}
    \item $x_{i,k}$: cantidad de producto $k$ producido en la planta $i$ $(\geq 0)$.
    \item $y_{i,j,k}$: cantidad de producto $k$ enviada desde $i$ a $j$ $(\geq 0)$
    \item $z_{i,k}$: Indica si se produce el producto $k$ en la planta $i$, de modo que:
\end{itemize}

\[
z_{i,k} =
\begin{cases}
1, & \text{si se produce } k \text{ en la planta } i \\
0, & \text{otro caso}
\end{cases}
\]

Los parámetros del modelo estan descritas en el enunciado. La función objetivo será:

\[
\min \sum_{i=1}^{I} \sum_{k=1}^{K} \big(V_{i,k} \cdot x_{i,k} +  f_{i,k} \cdot z_{i,k} \big)
+ \sum_{i=1}^{I} \sum_{j=1}^{J} \sum_{k=1}^{K} C_{i,j,k} \cdot y_{i,j,k}
\]

Observamos que el primer término contiene los costos variables y fijos de producción de cada producto $k$ en la planta $i$. El segundo término corresponde a los costos de envío de cada producto $k$ desde la planta $i$ hacia la zona $j$. \\

Dado el contexto del problema y los parámetros entregados en el enunciado, las siguientes restricciones deben ser planteadas:

\begin{itemize}
    \item[i)] Satisfacción de demanda del producto $k$ en la zona $j$, desde las distintas plantas de producción:
    \begin{equation}
    \sum_{i=1}^{I} y_{i,j,k} = d_{j,k}
    \qquad \forall j,\ \forall k
    \label{ex:2:4:eq:14}
    \end{equation}
    \item[ii)] Cantidad de producción de $k$ en planta $i$ debe ser igual a lo enviado hacia las distintas zonas de consumo:
    \begin{equation}
        \sum_{j=1}^{J} y_{i,j,k} = x_{i,k} \qquad \forall i,\ \forall k
    \end{equation}
    \item[iii)] Respetar limites técnicos de cada planta $i$ por producción de $k$. Notar que se debe plantear utilizando $z_{i,k}$, similar a lo realizado en P3 con los despachos por generador, y en P2 con las potencia de carga/descarga de batería.
    
    \begin{equation}
        x_{i,k} \leq M_{i,k} \cdot z_{i,k}
    \qquad \forall k,\ \forall i
    \end{equation}

    \begin{equation}
        x_{i,k} \geq m_{i,k} \cdot z_{i,k}
    \qquad \forall k,\ \forall i
    \end{equation}

    \textbf{Observación:} La variable $z_{i,k}$ se utiliza en la restricción de limite y no en la función objetivo para evitar que la formulación sea no lineal.

    \item[iv)] Respetar límite de producción total de cada planta $i$:
    \begin{equation}
        \sum_{k=1}^K x_{i,k} \leq Q_i
    \qquad \forall i
    \end{equation}
\end{itemize}

De modo que la formulación será:

\[
\begin{aligned}
\min \quad
& \sum_{i=1}^{I} \sum_{k=1}^{K} \big(V_{i,k} \cdot x_{i,k} +  f_{i,k} \cdot z_{i,k} \big)
+ \sum_{i=1}^{I} \sum_{j=1}^{J} \sum_{k=1}^{K} C_{i,j,k} \cdot y_{i,j,k} \\
\text{s.a} \quad
& \sum_{i=1}^{I} y_{i,j,k} = d_{j,k} && \forall j,\ \forall k \\
& \sum_{j=1}^{J} y_{i,j,k} = x_{i,k} && \forall i,\ \forall k \\
& x_{i,k} \leq M_{i,k} \cdot z_{i,k} && \forall k,\ \forall i \\
& x_{i,k} \geq m_{i,k} \cdot z_{i,k} && \forall k,\ \forall i \\
& \sum_{k=1}^K x_{i,k} \leq Q_i && \forall i \\
& x_{i,k} \geq 0 && \forall i,\ \forall k \\
& y_{i,j,k} \geq 0 && \forall i,\forall j ,\ \forall k \\
& z_{i,k} \in \{0,1\} && \forall i,\ \forall k
\end{aligned}
\]

\textbf{Restricciones adicionales}

\begin{enumerate}
\item
\begin{equation}
    \sum_{k=1}^{K} z_{i,k} \leq K_1
\qquad \forall i
\end{equation}

\item
\begin{equation}
    \sum_{i=1}^{I} z_{i,k} \leq I_1
\qquad \forall k
\end{equation}

\item
\begin{equation}
    z_{3,k_0} \geq 1 - (z_{1,k_0} + z_{2,k_0})
\end{equation}

\item
Se define una nueva variable binaria:

\begin{equation}
    \omega_{i,j} =
\begin{cases}
1, & \text{si la planta } i \text{ abastece la zona } j \\
0, & \text{otro caso}
\end{cases}
\end{equation}

La variable $\omega_{i,j}$ se implementa para limitar el flujo de productos de $i$ a $j$, mediente la siguiente restricción:

\begin{equation}
    y_{i,j,k} \leq d_{j,k} \cdot \omega_{i,j}
\qquad \forall i,\ \forall j,\ \forall k
\label{ex:2:4:eq:21}
\end{equation}

Luego, para asegurar que una zona $j$ sea abastecida por una única planta $i$, se impone la siguiente restricción:

\begin{equation}
    \sum_{i=1}^{I} \omega_{i,j} = 1
\qquad \forall j
\label{ex:2:4:eq:22}
\end{equation}

Notar que mediante las ecuaciones \eqref{ex:2:4:eq:14}, \eqref{ex:2:4:eq:21} y \eqref{ex:2:4:eq:22} se garantiza que para cada zona $j$, solo existirá una planta $i$ desde la cual las cantidades enviadas $y_{i,j,k}$ serán distintas de $0$ para todos los productos.
\end{enumerate}
